\documentclass{situationist-preprint}

\title{Cognitive Enhancement and Fairness}
\msid{ETH-22-b}
\msfamily{ETH-22}
\msdiscipline{Normative and Applied Ethics}
\mstype{Companion manuscript}
\mskeywords{cognitive enhancement; fairness; positional goods; access; examinations; distributive justice}

\begin{document}
\maketitle

\begin{abstract}
Cognitive enhancement is often thought to raise distinctive questions of fairness: students using stimulants are said to cheat, and unequal access to enhancers to entrench advantage. We apply the reply of ETH-22-a to the objection from positional competition and argue that these questions fall under the general fairness constraint on medical interventions and require no further one. Where the benefits of a cognitive enhancement are positional, as in competitive examinations, fairness may require regulation by the rules of the competition; where they are not, as for memory and learning, it requires only fair access. Neither verdict depends on the intervention's being an enhancement rather than a therapy.
\end{abstract}

\section{The argument}

ETH-22-a holds that enhancement is subject to the constraints of consent and fairness that govern every medical intervention, and to no further constraint (its Theorem~3.3). Its Principle~P2 requires that access be distributed justly and that effects on third parties, including effects on their competitive position, be consistent with fair terms of cooperation. Cognitive enhancement, by pharmacological or other means, raises both issues \citep{Greely2008,BostromSandberg2009}.

\begin{definition}\label{d:positional}
A benefit is \emph{positional} in a domain if its value to its bearer in that domain depends on its being greater than the corresponding benefits of others.
\end{definition}

\begin{proposition}\label{p:competition}
Where the benefit of a cognitive enhancement is positional in a competition whose rules are fair, the competition may permissibly restrict its use; where it is not positional, fairness gives no reason to restrict it beyond the requirement of fair access.
\end{proposition}

\begin{proof}[Argument]
A competitive examination ranks candidates for a limited good. If an enhancer raises the scores of those who take it, its benefit in that domain is positional, and its unrestricted use either confers an advantage on users at others' expense or, if general, leaves rankings as they were while imposing costs and risks on all. Either effect bears on the competitive position of third parties, which Principle~P2 of ETH-22-a covers; the competition's rules may therefore restrict the enhancer, as sporting bodies restrict performance-enhancing drugs. Better memory in old age, faster learning of a language or improved attention at work benefit their bearers whether or not others have them. These effects on others are not losses of position, and Principle~P2 then requires only that access be just, which is the requirement for any medical good.
\end{proof}

The proposition applies the second clause of the fairness principle without any premise about enhancement as such. A therapeutic treatment for attention deficit raises the same question in an examination, and many examining bodies already settle it by rules about accommodations. The difference between the cases is a difference in the effects of the intervention on others, not in its relation to normal functioning, which by Lemma~3.2 of ETH-22-a has no moral authority of its own.

\begin{objection}
Even non-positional cognitive enhancements will widen inequality, since the well-off will obtain them first.
\end{objection}

\begin{reply}
This is the question of access, and it arises for every valuable medical good, including vaccines and therapies that are at first scarce and expensive. It is answered by the principles of distributive justice that govern access to medical goods generally \citep{Buchanan2011}. It gives no reason to forbid the enhancement, since forbidding it secures equality only by denying the benefit to everyone.
\end{reply}

\section{Conclusion}

The fairness questions raised by cognitive enhancement are questions about positional effects and about access. Both are governed by the fairness constraint that applies to all medical interventions, and the answer to each depends on the domain and the distribution, not on whether the intervention enhances.

\begin{thebibliography}{9}\small
\bibitem[Bostrom and Sandberg(2009)]{BostromSandberg2009} Bostrom, N. and Sandberg, A. (2009). Cognitive enhancement: methods, ethics, regulatory challenges. \emph{Science and Engineering Ethics}, 15(3), 311--341.
\bibitem[Buchanan(2011)]{Buchanan2011} Buchanan, A. (2011). \emph{Beyond Humanity? The Ethics of Biomedical Enhancement}. Oxford University Press.
\bibitem[Greely et~al.(2008)]{Greely2008} Greely, H., Sahakian, B., Harris, J., Kessler, R.~C., Gazzaniga, M., Campbell, P. and Farah, M.~J. (2008). Towards responsible use of cognitive-enhancing drugs by the healthy. \emph{Nature}, 456(7223), 702--705.
\end{thebibliography}

\end{document}
